% GENERATED FILE. Edit canonical lessons, then run npm run generate:books.
% canonical-source-hash: fnv1a64:ed9e32a32cdebd49
% canonical-lessons: FA-C57-khers, FA-C57-shotor, FA-C57-khar, FA-C57-parvane, FA-C57-murche

\chapter{Bear, Camel, Donkey, Butterfly, Ant}
\label{ch:fa-bear-camel-donkey-butterfly-ant}
\hlchaptermodality{57}

\begin{chapteropening}
\textbf{By the end of this chapter:} \emph{I can say a bear, a camel, a donkey, a butterfly and an ant.}

Complete the last lesson of chapter 57: I can say a bear, a camel, a donkey, a butterfly and an ant.
\end{chapteropening}

\section[\texorpdfstring{\fa{خرس} (khers)}{خرس (khers)}]{\fa{خرس} (khers) — a bear}

\label{lesson:FA-C57-khers}



\begin{quote}
Before the new one: say the Persian for a monkey, then the Persian for a snake.
\end{quote}

\subsection*{You'll want to know: \fa{خرس}}
\textbf{\fa{خرس}} — \emph{khers} — \textquotedblleft{}a bear\textquotedblright{}.

\begin{cousinweb}
\emph{khers} is a native Persian word, inherited rather than borrowed; it is usually traced to Middle Persian \emph{xirs}, \textquotedblleft{}bear\textquotedblright{}. EtymDB takes it back to the Indo-European word for \textquotedblleft{}bear\textquotedblright{} behind Latin \emph{ursus} and Greek \emph{arktos}.
\end{cousinweb}

\begin{grammarlens}[title={What you've built}]
One more word for this chapter.
\end{grammarlens}

\subsection*{Your turn}
\begin{itemize}
\raggedright
  \item \emph{Say it:} \emph{khers}
  \item \emph{Say it:} \emph{khers}, once more
  \item \emph{Say it:} say \emph{khers}
  \item \emph{Recall:} say the Persian for a monkey, then the Persian for a snake, then say \emph{khers} again
\end{itemize}

\begin{tcolorbox}[breakable,colback=teal!4,colframe=teal!35!black,title={Before you move on}]
What does \fa{خرس} mean? (\textquotedblleft{}A bear\textquotedblright{}.) Say it once more.
\end{tcolorbox}
\section[\texorpdfstring{\fa{شتر} (shotor)}{شتر (shotor)}]{\fa{شتر} (shotor) — a camel}

\label{lesson:FA-C57-shotor}



\begin{quote}
Before the new one: say the Persian for a snake, then the Persian for a bear.
\end{quote}

\subsection*{You'll want to know: \fa{شتر}}
\textbf{\fa{شتر}} — \emph{shotor} — \textquotedblleft{}a camel\textquotedblright{}.

\begin{cousinweb}
\emph{shotor} is a native Persian word, inherited rather than borrowed: EtymDB traces it to Proto-Iranian \emph{*uštrah}, \textquotedblleft{}camel\textquotedblright{}, and Middle Persian had \emph{uštar}.
\end{cousinweb}

\begin{grammarlens}[title={What you've built}]
One more word for this chapter.
\end{grammarlens}

\subsection*{Your turn}
\begin{itemize}
\raggedright
  \item \emph{Say it:} \emph{shotor}
  \item \emph{Say it:} \emph{shotor}, once more
  \item \emph{Say it:} say \emph{shotor}
  \item \emph{Recall:} say the Persian for a snake, then the Persian for a bear, then say \emph{shotor} again
\end{itemize}

\begin{tcolorbox}[breakable,colback=teal!4,colframe=teal!35!black,title={Before you move on}]
What does \fa{شتر} mean? (\textquotedblleft{}A camel\textquotedblright{}.) Say it once more.
\end{tcolorbox}
\section[\texorpdfstring{\fa{خر} (khar)}{خر (khar)}]{\fa{خر} (khar) — a donkey}

\label{lesson:FA-C57-khar}



\begin{quote}
Before the new one: say the Persian for a bear, then the Persian for a camel.
\end{quote}

\subsection*{You'll want to know: \fa{خر}}
\textbf{\fa{خر}} — \emph{khar} — \textquotedblleft{}a donkey\textquotedblright{}.

\begin{cousinweb}
\emph{khar} is a native Persian word, inherited rather than borrowed; it is usually traced to Middle Persian \emph{xar}, \textquotedblleft{}donkey\textquotedblright{}.
\end{cousinweb}

\begin{grammarlens}[title={What you've built}]
One more word for this chapter.
\end{grammarlens}

\subsection*{Your turn}
\begin{itemize}
\raggedright
  \item \emph{Say it:} \emph{khar}
  \item \emph{Say it:} \emph{khar}, once more
  \item \emph{Say it:} say \emph{khar}
  \item \emph{Recall:} say the Persian for a bear, then the Persian for a camel, then say \emph{khar} again
\end{itemize}

\begin{tcolorbox}[breakable,colback=teal!4,colframe=teal!35!black,title={Before you move on}]
What does \fa{خر} mean? (\textquotedblleft{}A donkey\textquotedblright{}.) Say it once more.
\end{tcolorbox}
\section[\texorpdfstring{\fa{پروانه} (parvâne)}{پروانه (parvâne)}]{\fa{پروانه} (parvâne) — a butterfly}

\label{lesson:FA-C57-parvane}



\begin{quote}
Before the new one: say the Persian for a camel, then the Persian for a donkey.
\end{quote}

\subsection*{You'll want to know: \fa{پروانه}}
\textbf{\fa{پروانه}} — \emph{parvâne} — \textquotedblleft{}a butterfly\textquotedblright{}.

\begin{cousinweb}
\emph{parvâne} continues Middle Persian \emph{parwānag} (EtymDB), which meant \textquotedblleft{}guide, leader\textquotedblright{}. How the word came to name the butterfly and the moth is not settled.
\end{cousinweb}

\begin{grammarlens}[title={What you've built}]
One more word for this chapter.
\end{grammarlens}

\subsection*{Your turn}
\begin{itemize}
\raggedright
  \item \emph{Say it:} \emph{parvâne}
  \item \emph{Say it:} \emph{parvâne}, once more
  \item \emph{Say it:} say \emph{parvâne}
  \item \emph{Recall:} say the Persian for a camel, then the Persian for a donkey, then say \emph{parvâne} again
\end{itemize}

\begin{tcolorbox}[breakable,colback=teal!4,colframe=teal!35!black,title={Before you move on}]
What does \fa{پروانه} mean? (\textquotedblleft{}A butterfly\textquotedblright{}.) Say it once more.
\end{tcolorbox}
\section[\texorpdfstring{\fa{مورچه} (murche)}{مورچه (murche)}]{\fa{مورچه} (murche) — an ant}

\label{lesson:FA-C57-murche}



\begin{quote}
Before the new one: say the Persian for a donkey, then the Persian for a butterfly.
\end{quote}

\subsection*{You'll want to know: \fa{مورچه}}
\textbf{\fa{مورچه}} — \emph{murche} — \textquotedblleft{}an ant\textquotedblright{}.

\begin{cousinweb}
\emph{murche} is \emph{mur}, \textquotedblleft{}ant\textquotedblright{}, with the small-thing ending \emph{-che}. \emph{mur} is a native word, usually traced to Middle Persian \emph{mōr}.
\end{cousinweb}

\begin{grammarlens}[title={What you've built}]
One more word for this chapter.
\end{grammarlens}

\subsection*{Your turn}
\begin{itemize}
\raggedright
  \item \emph{Say it:} \emph{murche}
  \item \emph{Say it:} \emph{murche}, once more
  \item \emph{Say it:} say \emph{murche}
  \item \emph{Recall:} say the Persian for a donkey, then the Persian for a butterfly, then say \emph{murche} again
\end{itemize}

\begin{tcolorbox}[breakable,colback=teal!4,colframe=teal!35!black,title={Before you move on}]
What does \fa{مورچه} mean? (\textquotedblleft{}An ant\textquotedblright{}.) Say it once more.
\end{tcolorbox}
